\hypertarget{bitwise-operations}{%
\chapter{BITWISE OPERATIONS}\label{bitwise-operations}}

\hypertarget{bitwise-operators}{%
\section{6.1 Bitwise Operators}\label{bitwise-operators}}

\begin{lstlisting}[language={C++}]
#include <iostream>
using namespace std;

int main() {
    unsigned char a = 5;   // 0101
    unsigned char b = 3;   // 0011
    
    // AND
    cout << (a & b) << endl;  // 0001 = 1
    
    // OR
    cout << (a | b) << endl;  // 0111 = 7
    
    // XOR
    cout << (a ^ b) << endl;  // 0110 = 6
    
    // NOT (bitwise complement)
    cout << (~a) << endl;     // 1010 = 250 (for unsigned char)
    
    // Left shift
    cout << (a << 1) << endl; // 1010 = 10
    
    // Right shift
    cout << (b >> 1) << endl; // 0001 = 1
    
    return 0;
}
\end{lstlisting}

\hypertarget{bit-manipulation-techniques}{%
\section{6.2 Bit Manipulation
Techniques}\label{bit-manipulation-techniques}}

\begin{lstlisting}[language={C++}]
#include <iostream>
using namespace std;

int main() {
    unsigned int num = 5;  // 0101
    
    // Check if bit is set
    int bit_pos = 2;
    bool is_set = (num >> bit_pos) & 1;
    cout << "Bit " << bit_pos << " is: " << is_set << endl;
    
    // Set a bit
    num |= (1 << 1);  // Set bit 1
    cout << "After setting bit 1: " << num << endl;  // 7 (0111)
    
    // Clear a bit
    num &= ~(1 << 1);  // Clear bit 1
    cout << "After clearing bit 1: " << num << endl;  // 5 (0101)
    
    // Toggle a bit
    num ^= (1 << 0);  // Toggle bit 0
    cout << "After toggling bit 0: " << num << endl;  // 4 (0100)
    
    // Count set bits
    unsigned int count = 0;
    unsigned int temp = num;
    while (temp) {
        count += temp & 1;
        temp >>= 1;
    }
    cout << "Number of set bits: " << count << endl;
    
    return 0;
}
\end{lstlisting}

\begin{center}\rule{0.5\linewidth}{0.5pt}\end{center}
